% ============================================================================
% Omar Khayyam - Risala fi al-Jabr wa-al-Muqabala
% Arabic Translation of the Treatise on Algebra
% ============================================================================
\documentclass[11pt,a4paper]{book}
\usepackage{fontspec}
\usepackage{amsmath,amssymb}
\usepackage{geometry}
\usepackage{xcolor}
\usepackage{titlesec}
\usepackage{fancyhdr}
\usepackage{hyperref}
\usepackage{tikz}
\usepackage{booktabs}
\usepackage{longtable}
\usepackage{enumitem}

\usetikzlibrary{arrows.meta,calc,intersections}

\geometry{paperwidth=210mm,paperheight=297mm,left=25mm,right=25mm,top=30mm,bottom=28mm,headheight=14pt}

\setmainfont{Arial}[BoldFont=Arial,ItalicFont=Arial,BoldItalicFont=Arial]\setsansfont{Arial}[BoldFont=Arial,ItalicFont=Arial,BoldItalicFont=Arial]% Arabic font setup with RTL support
\TeXXeTstate=1
\newfontfamily\arabicfont[Script=Arabic,BoldFont=Arial,ItalicFont=Arial,BoldItalicFont=Arial]{Arial}\newfontfamily\arabicfontbold[Script=Arabic,BoldFont=Arial,ItalicFont=Arial,BoldItalicFont=Arial]{Arial}\newcommand{\ar}[1]{\leavevmode\arabicfont\beginR #1\endR}
\newcommand{\arBold}[1]{\leavevmode\arabicfontbold\beginR #1\endR}
\newcommand{\arLarge}[1]{{\Large\leavevmode\arabicfont\beginR #1\endR}}
\newcommand{\arHuge}[1]{{\Huge\leavevmode\arabicfont\beginR #1\endR}}
\newcommand{\arLARGE}[1]{{\LARGE\leavevmode\arabicfont\beginR #1\endR}}
\newcommand{\arLargebf}[1]{{\Large\bfseries\leavevmode\arabicfontbold\beginR #1\endR}}

% Colors
\definecolor{titlecolor}{RGB}{139,69,19}
\definecolor{sectioncolor}{RGB}{101,67,33}
\definecolor{rulecolor}{RGB}{160,82,45}
\definecolor{boxbg}{RGB}{255,250,240}
\definecolor{boxborder}{RGB}{139,69,19}
\definecolor{shadecolor}{RGB}{255,250,240}
\renewenvironment{leftbar}{\def\FrameCommand{\color{rulecolor}\vrule width 1.5pt\hspace{10pt}}\MakeFramed{\advance\hsize-\width\FrameRestore}}{\endMakeFramed}

% Arabic section formatting
\titleformat{\section}
{\color{sectioncolor}\Large\bfseries\arabicfontbold}
{\thesection}{1em}{\beginR}
\titleformat{\subsection}
{\color{sectioncolor}\large\bfseries\arabicfontbold}
{\thesubsection}{1em}{\beginR}
\titleformat{\subsubsection}
{\normalsize\bfseries\arabicfontbold}
{\thesubsubsection}{1em}{\beginR}

% Header/footer
\pagestyle{fancy}
\fancyhf{}
\fancyhead[R]{\arabicfont\beginR صناعة الجبر والمقابلة\endR}
\fancyhead[L]{\arabicfont\beginR عمر الخيام\endR}
\fancyfoot[C]{\thepage}
\renewcommand{\headrulewidth}{0.4pt}

% Custom box for propositions (using framed package instead of tcolorbox)
\usepackage{framed}
\usepackage{xcolor}

% Proposition box with colored background
\newenvironment{propbox}[1][]
  {\def\propboxtitle{#1}%
   \begin{snugshade}%
   \ifx\propboxtitle\empty\else
     \noindent{\bfseries\arabicfontbold\beginR #1\endR}\par\vspace{2pt}
   \fi}
  {\end{snugshade}}

% Definition box with left rule
\newenvironment{defbox}
  {\begin{leftbar}}
  {\end{leftbar}}

\newcommand{\vs}{\vspace{0.5em}}
\begin{document}

% ============================================================================
% TITLE PAGE
% ============================================================================
\begin{titlepage}
\begin{center}

\vspace*{1.5cm}

{\color{titlecolor}\rule{\textwidth}{1.5pt}}

\vspace{1cm}

{\arHuge{عمر الخيام}}

\vspace{0.5cm}

{\large\itshape (1048--1131 م / 440--517 هـ)}

\vspace{1.5cm}

{\color{titlecolor}\rule{0.7\textwidth}{0.8pt}}

\vspace{1.5cm}

{\arHuge{رسالة في براهين الجبر والمقابلة}}

\vspace{0.8cm}

{\arLARGE{صناعة الجبر والمقابلة}}

\vspace{1cm}

{\arLarge{ترجمة عربية مع الحفاظ على الرموز الرياضية الحديثة}}

\vspace{0.5cm}

{\normalsize\textit{Arabic Translation with Modern Mathematical Notation}}

\vspace{1.5cm}

{\color{titlecolor}\rule{0.7\textwidth}{0.8pt}}

\vspace{1cm}

{\arLarge{يشتمل على تصنيف المعادلات التكعيبية والحلول الهندسية باستخدام القطوع المخروطية}}

\vspace{2cm}

{\small\textit{Based on the Lahore Manuscript (13th Century)}}

\vspace{0.3cm}

{\small\textit{(Smith Oriental MS 45, Columbia University Rare Book Library)}}

\vspace{1cm}

{\color{titlecolor}\rule{\textwidth}{1.5pt}}

\end{center}
\end{titlepage}

% ============================================================================
% COPYRIGHT PAGE
% ============================================================================
\thispagestyle{empty}
\vspace*{8cm}
\begin{center}
{\small\ar{هذا الكتاب يقدم نسخة مصنفة من عمل عمر الخيام في الجبر،}}

\vspace{0.3cm}

{\small\ar{وهو من أبرز الإنجازات في الرياضيات الإسلامية الوسطى.}}

\vspace{0.5cm}

{\small\ar{يتضمن المعالجة الأولى المنهجية للمعادلات التكعيبية،}}

{\small\ar{حيث تم حل جميع الأصناف الخمسة والعشرين باستخدام القطوع المخروطية.}}

\vspace{1cm}

{\footnotesize\textit{Typeset in \LaTeX\ with Arabic support via fontspec and XeLaTeX.}}

{\footnotesize\textit{Arabic font: Noto Naskh Arabic}}

\vspace{0.5cm}

{\footnotesize\textit{This edition prepared for scholarly and educational purposes.}}

\end{center}
\clearpage

% ============================================================================
% TABLE OF CONTENTS
% ============================================================================
\tableofcontents
\clearpage

% ============================================================================
% المقدمة
% ============================================================================
\section{\ar{المقدمة}}

\ar{عمر الخيام}، اسمه الكامل \ar{غياث الدين أبو الفتح عمر بن إبراهيم الخيامي النيسابوري}، من أبرز علماء الرياضيات والفلك والشعر في العالم الإسلامي الوسيط. وُلد في نيسابور عام 1048 م، وقضى معظم حياته في سمرقند وأصفهان، حيث أنتج أعمالاً رائدة في الرياضيات والفلك.

تُعد \textbf{\ar{رسالة في براهين الجبر والمقابلة}} (\textit{صناعة الجبر والمقابلة}) ذروة الجبر الإسلامي المبكر. في هذا العمل، يقوم الخيام بما يلي:

\begin{itemize}[noitemsep]
\item \ar{تقديم تعريفات دقيقة للكميات الأساسية في الجبر}
\item \ar{تصنيف جميع المعادلات حتى الدرجة الثالثة إلى خمسة وعشرين صنفاً}
\item \ar{حل أربعة عشر معادلة تكعيبية غير قابلة للاختزال باستخدام التراكيب الهندسية بالقطوع المخروطية}
\item \ar{التمييز بين الحلول الحسابية والحلول الهندسية}
\end{itemize}

هذه الرسالة، التي ألفها حوالي عام 1079 م، أثرت بعمق في الرياضيين اللاحقين بما فيهم \ar{شرف الدين الطوسي}، ومن خلال الترجمات إلى اللغات الأوروبية، ساهمت في تطوير الجبر في الغرب.

% ============================================================================
% مقدمة المؤلف
% ============================================================================
\subsection{\ar{مقدمة المؤلف}}

يفتح الخيام رسالته بمقدمة يعترف فيها بإنجازات سابقيه ومحدوديتهم:

\begin{propbox}[\arBold{مقدمة المؤلف}]

\ar{الحمد لله رب العالمين، والعاقبة للمتقين، ولا عدوان إلا على الظالمين، وصلى الله على سيدنا محمد وآله أجمعين.}

\vspace{0.3cm}

\textit{الحمد لله رب العالمين، والعاقبة للمتقين، ولا عدوان إلا على الظالمين، وصلِّ الله على سيدنا محمد وآله أجمعين.}

\vspace{0.5cm}

\ar{أما بعد، فإني قد سبقني في هذا الفن جماعة من أهل العلم، قدموا فيه أصناف المسائل، وبينوا فيه أصول المسائل، وبرهنوا على بعضها، ولم يتفق لهم برهان على البعض، ولم يفرقوا فيها بين ما يمكن أن يوجد بالقطوع المخروطية، وبين ما لا يمكن أن يوجد إلا بأساليب هندسية أخرى. وقد أملت أن أجمع هذا العلم في رسالة، وأبين فيها أصناف المسائل على وجه التمييز، وأضع لكل صنف برهاناً صحيحاً، وذلك أمر لم يسبقني إليه أحد من المتقدمين.}

\vspace{0.3cm}

\textit{الآن: في هذا الفن سبقني مجموعة من أهل العلم، قدموا فيه أصناف المسائل، وبينوا فيها أصول المسائل، وبرهنوا على بعضها، ولكن لم ينجحوا في إثباتها جميعاً؛ ولم يميزوا بين ما يمكن إيجاده بالقطوع المخروطية وما لا يمكن إيجاده إلا بطرق هندسية أخرى. لقد رغبت في جمع هذا العلم في رسالة، وتوضيح أصناف المسائل مع التمييز الصحيح، ووضع برهان صحيح لكل صنف --- وهذا شيء لم يسبقني إليه أحد من السابقين.}

\end{propbox}

% ============================================================================
% في طبيعة هذه الرسالة
% ============================================================================
\subsection{\ar{في طبيعة هذه الرسالة}}

\ar{يُعرف الخيام الجبر بأنه ``علم المعادلات'' (علم المعادلات). ويميز بعناية بين:}

\begin{enumerate}[noitemsep]
\item \textbf{\ar{المعادلات البسيطة:}} \ar{التي يمكن اختزالها إلى أشكال خطية أو تربيعية، وتحل بالطرق الحسابية أو الإقليدية}
\item \textbf{\ar{المعادلات المركبة (غير القابلة للاختزال):}} \ar{المعادلات التكعيبية التي تتطلب تراكيب القطوع المخروطية}
\end{enumerate}

\begin{propbox}[AutoFakeBold=true,AutoFakeSlant=true]
\ar{فنقول: إن الجبر والمقابلة علم يُستخرج به المجهول من المعلوم المفروض، إذا كان بينهما نسبة تقتضي ذلك. وهذا العلم يتضمن الأعداد والجذور والأموال والكعاب.}

\vspace{0.3cm}

\textit{نقول: الجبر والمقابلة علم يُستخرج به المجهول من المعطى المعروف، عندما تكون هناك علاقة تتطلب ذلك. وهذا العلم يشمل الأعداد والجذور والأموال والكعاب.}
\end{propbox}

\clearpage

% ============================================================================
% تعريفات الكميات الأساسية
% ============================================================================
\section{\ar{تعريفات الكميات الأساسية}}

\ar{يبدأ الخيام بتعريفات دقيقة للكميات الجبرية، مستمراً وموسعاً لتقاليد الخوارزمي:}

\begin{defbox}
\ar{أصل هذا العلم: أن العدد المفرد هو ما لا ينسب إلى جذر، ولا إلى مال، ولا إلى كعب، ولا إلى ما يتركب منها.}

\vspace{0.3cm}

\textit{أساس هذا العلم هو: العدد البسيط هو ما لا ينسب إلى جذر ولا إلى مال ولا إلى كعب ولا إلى ما يُركب منها.}
\end{defbox}

% ============================================================================
% الجذر
% ============================================================================
\subsection{\ar{الجذر (الشيء)}}

\ar{والجذر هو كل شيء مضروب في نفسه من الواحد وما فوقه من الأعداد، وما دونه من الكسور. فإذا ضرب في نفسه فهو مال، وإذا ضرب المال في الجذر فهو كعب، وإذا ضرب الكعب في الجذر فهو مال مال، وإذا ضرب مال المال في الجذر فهو مال كعب، وإذا ضرب مال الكعب في الجذر فهو كعب كعب.}

\vspace{0.3cm}

\textit{الجذر هو كل كمية مضروبة في نفسها، بدءاً من الواحد وأعلى في الأعداد الصحيحة، وأقل في الكسور. إذا ضُرب في نفسه فهو مال؛ وإذا ضُرب المال في الجذر فهو كعب؛ وإذا ضُرب الكعب في الجذر فهو مال مال؛ وإذا ضُرب مال المال في الجذر فهو مال كعب؛ وإذا ضُرب مال الكعب في الجذر فهو كعب كعب.}

% ============================================================================
% المال
% ============================================================================
\subsection{\ar{المال (المربع)}}

\ar{والمال هو ما ينتج من ضرب الجذر في نفسه. والكعب هو ما ينتج من ضرب المال في الجذر. وما فوق الكعب فهو مركب من هذه الأجناس.}

\vspace{0.3cm}

\textit{المال هو ما ينتج من ضرب الجذر في نفسه. الكعب هو ما ينتج من ضرب المال في الجذر. وما فوق الكعب فهو مركب من هذه الأجناس.}

% ============================================================================
% تصنيف الأجناس الجبرية
% ============================================================================
\subsection{\ar{تصنيف الأجناس الجبرية}}

\ar{يعدّ الخيام الأجناس البسيطة والمركبة بشكل منهجي:}

\begin{center}
\renewcommand{\arraystretch}{1.4}
\begin{tabular}{ccc}
\toprule
\arBold{الحدود الجبرية} & \textbf{الاسم اللاتيني} & \textbf{الرموز الحديثة} \\
\midrule
\ar{جذر} & \textit{jidhr} / radix & $x$ \\
\ar{مال} & \textit{māl} & $x^2$ \\
\ar{كعب} & \textit{kaʿb} & $x^3$ \\
\ar{مال مال} & \textit{māl māl} & $x^4$ \\
\ar{مال كعب} & \textit{māl kaʿb} & $x^5$ \\
\ar{كعب كعب} & \textit{kaʿb kaʿb} & $x^6$ \\
\bottomrule
\end{tabular}
\end{center}

\vspace{0.3cm}

\ar{في هذا التصنيف، نرى بوضوح التقدم الذي أحرزه الخيام عن سابقيه. فبينما اقتصرت أعمال الخوارزمي على الأجناس حتى الكعب، توسع الخيام في الأجناس المركبة وأعطى لها أسماء واضحة ومنهجية.}

\clearpage

% ============================================================================
% تصنيف المعادلات
% ============================================================================
\section{\ar{تصنيف المعادلات}}

\ar{يشكل هذا القسم قلب رسالة الخيام. يصنف جميع المعادلات حتى الدرجة الثالثة إلى \textbf{خمسة وعشرين صنفاً}. وبالاتباع لعوائد عصره، تكون جميع المعاملات موجبة، فتُعتبر معادلات مثل $x^3 + cx = bx^2 + d$ و $x^3 + bx^2 = cx + d$ من أصناف مختلفة.}

% ============================================================================
% المعادلات القابلة للحل بالطرق الحسابية أو الإقليدية
% ============================================================================
\subsection{\ar{المعادلات القابلة للحل بالطرق الحسابية أو الإقليدية}}

\subsubsection{\ar{المعادلات البسيطة (الأصناف 1--6)}}

\ar{تتكون المجموعة الأولى من المعادلات الخطية والتربيعية، التي تحل بدون القطوع المخروطية:}

\begin{enumerate}[label=\arabic*.,noitemsep]
\item \ar{جذور تساوي أعداداً} --- \ar{الجذور تساوي الأعداد}: $bx = c$
\item \ar{أموال تساوي جذوراً} --- \ar{الأموال تساوي الجذور}: $x^2 = bx$
\item \ar{أموال تساوي أعداداً} --- \ar{الأموال تساوي الأعداد}: $x^2 = c$
\item \ar{أموال وجذور تساوي أعداداً} --- \ar{الأموال والجذور تساوي الأعداد}: $x^2 + bx = c$
\item \ar{أموال وأعداد تساوي جذوراً} --- \ar{الأموال والأعداد تساوي الجذور}: $x^2 + c = bx$
\item \ar{جذور وأعداد تساوي أموالاً} --- \ar{الجذور والأعداد تساوي الأموال}: $bx + c = x^2$
\end{enumerate}

\ar{كانت هذه الأصناف الستة معروفة منذ الخوارزمي وتحل بالطرق الابتدائية.}

% ============================================================================
% المعادلات التكعيبية الحدية
% ============================================================================
\subsubsection{\ar{المعادلات التكعيبية الحدية (الأصناف 7--12)}}

\ar{وأما الكعاب والأموال والجذور التي لا تختلف فيها الأجناس، فإنها تنقسم إلى اثني عشر نوعاً: منها ما يحل بالجبر، ومنها ما يحل بالقطوع المخروطية.}

\begin{enumerate}[label=\arabic*.,resume,noitemsep]
\item \ar{كعب يساوي مالاً} --- \ar{الكعب يساوي المال}: $x^3 = bx^2$
\item \ar{كعب يساوي جذراً} --- \ar{الكعب يساوي الجذر}: $x^3 = cx$
\item \ar{كعب يساوي عدداً} --- \ar{الكعب يساوي العدد}: $x^3 = d$
\item \ar{مال يساوي كعباً} --- \ar{المال يساوي الكعب}: $bx^2 = x^3$
\item \ar{جذر يساوي كعباً} --- \ar{الجذر يساوي الكعب}: $cx = x^3$
\item \ar{عدد يساوي كعباً} --- \ar{العدد يساوي الكعب}: $d = x^3$
\end{enumerate}

\ar{هذه المعادلات الحدية تختزل إلى أشكال أبسط بالقسمة.}

% ============================================================================
% المعادلات الثلاثية القابلة للاختزال إلى تربيعيات
% ============================================================================
\subsubsection{\ar{المعادلات الثلاثية القابلة للاختزال (الأصناف 13--16)}}

\begin{enumerate}[label=\arabic*.,resume,noitemsep]
\item \ar{كعاب وأموال تساوي جذوراً} --- \ar{الكعاب والأموال تساوي الجذور}: $x^3 + bx^2 = cx$
\item \ar{كعاب وجذور تساوي أموالاً} --- \ar{الكعاب والجذور تساوي الأموال}: $x^3 + cx = bx^2$
\item \ar{أموال وجذور تساوي كعاباً} --- \ar{الأموال والجذور تساوي الكعاب}: $bx^2 + cx = x^3$
\item \ar{كعاب وأموال وأعداد} --- \ar{أشكال خاصة قابلة للاختزال}
\end{enumerate}

\clearpage

% ============================================================================
% المعادلات التكعيبية غير القابلة للاختزال
% ============================================================================
\subsection{\ar{المعادلات التكعيبية غير القابلة للاختزال (الأصناف الأربعة عشر)}}

\ar{يُعد إنجاز رسالة الخيام الأبرز حله المنهجي للـ\textbf{أربعة عشر معادلة تكعيبية غير قابلة للاختزال}. وهذه تتطلب تقاطع القطوع المخروطية لإيجاد حلها الهندسي. نقدمها هنا مع أسمائها العربية ورموزها الحديثة:}

\begin{propbox}[\arBold{الأصناف الأربعة عشر غير القابلة للاختزال}]

\textbf{\ar{المجموعة الأولى: معادلات ثلاثية ذات قوى متعاقبة}}

\begin{enumerate}[label=\arabic*.,resume,noitemsep]
\item \ar{كعب ومال يساويان عدداً} --- \ar{الكعب والمال يساويان العدد}:
$$x^3 + bx^2 = d$$
\item \ar{كعب وعدد يساويان مالاً} --- \ar{الكعب والعدد يساويان المال}:
$$x^3 + d = bx^2$$
\item \ar{كعب يساوي مالاً وعدداً} --- \ar{الكعب يساوي المال والعدد}:
$$x^3 = bx^2 + d$$
\end{enumerate}

\vspace{0.5cm}

\textbf{\ar{المجموعة الثانية: معادلات ثلاثية ذات قوى مختلطة}}

\begin{enumerate}[label=\arabic*.,resume,noitemsep]
\setcounter{enumi}{18}
\item \ar{كعب وجذور يساويان عدداً} --- \ar{الكعب والجذور يساويان العدد}:
$$x^3 + cx = d$$
\item \ar{كعب وعدد يساويان جذوراً} --- \ar{الكعب والعدد يساويان الجذور}:
$$x^3 + d = cx$$
\item \ar{كعب يساوي جذوراً وعدداً} --- \ar{الكعب يساوي الجذور والعدد}:
$$x^3 = cx + d$$
\end{enumerate}

\vspace{0.5cm}

\textbf{\ar{المجموعة الثالثة: معادلات رباعية --- ثلاثة حدود في طرف}}

\begin{enumerate}[label=\arabic*.,resume,noitemsep]
\setcounter{enumi}{21}
\item \ar{كعب وأموال يساويان جذوراً وعدداً} --- \ar{الكعب والأموال يساويان الجذور والعدد}:
$$x^3 + bx^2 = cx + d$$
\item \ar{كعب وجذور يساويان أموالاً وعدداً} --- \ar{الكعب والجذور يساويان الأموال والعدد}:
$$x^3 + cx = bx^2 + d$$
\item \ar{كعب وأعداد يساويان أموالاً وجذوراً} --- \ar{الكعب والأعداد يساويان الأموال والجذور}:
$$x^3 + d = bx^2 + cx$$
\end{enumerate}

\vspace{0.5cm}

\textbf{\ar{المجموعة الرابعة: معادلات رباعية --- حدان في كل طرف}}

\begin{enumerate}[label=\arabic*.,resume,noitemsep]
\setcounter{enumi}{24}
\item \ar{كعب يساوي أموالاً وجذوراً وأعداداً} --- \ar{الكعب يساوي الأموال والجذور والأعداد}:
$$x^3 = bx^2 + cx + d$$
\end{enumerate}

\end{propbox}

\clearpage

% ============================================================================
% في الحل الهندسي للمعادلات التكعيبية
% ============================================================================
\section{\ar{في الحل الهندسي للمعادلات التكعيبية}}

% ============================================================================
% المبرهنة الأساسية
% ============================================================================
\subsection{\ar{المبرهنة الأساسية (القضية التمهيدية)}}

\ar{قبل تقديم حلوله، يثبت الخيام مبرهنة أساسية عن الأجسام المستطيلة السطوح (المتوازيات المستطيلات):}

\begin{propbox}[\arBold{مقدمة في المجسمات المتساوية}]

\ar{إذا كان لنا مكعب مستطيل السطوح، قاعدته مربعة، وارتفاعه خط معلوم، وأردنا أن نبني على قاعدة مربعة أخرى مكعباً مستطيلاً مساوياً له، فنحن نجد خطاً يكون النسبة بين ضلع القاعدة الأولى وضلع القاعدة الثانية كالنسبة بين ذلك الخط والارتفاع.}

\vspace{0.3cm}

\textit{إذا كان لدينا متوازي مستطيلات قاعدته مربعة وارتفاعه خط معلوم، وأردنا أن نبني على قاعدة مربعة أخرى متوازياً مستطيلات مساوياً له، فإننا نجد خطاً بحيث تكون النسبة بين ضلع القاعدة الأولى وضلع القاعدة الثانية كنسبة ذلك الخط إلى الارتفاع.}

\end{propbox}

% ============================================================================
% حل كعب وجذور يساويان عدداً
% ============================================================================
\subsection{\ar{حل ``كعب وجذور يساويان عدداً'' ($x^3 + cx = d$)}}

\ar{هذا هو الحالة الأشهر في رسالة الخيام --- المعادلة $x^3 + cx = d$ (كعب وجذور يساويان عدداً). يحلها الخيام بتقاطع \textbf{القطع المكافئ} و\textbf{الدائرة}.}

\begin{propbox}[$x^3 + cx = d$]

\ar{فلنضرب مثالاً: كعب وجذور يساويان عدداً. فلنجعل العدد مكعباً مستطيل السطوح، قاعدته مربعة، أضلاعها مساوية لعدد الجذور، وارتفاعه خط معلوم. ثم نأخذ قطعاً مكافئاً رأسه نقطة، ومحوره خط مستقيم، ومعلمه مساوٍ لضلع القاعدة. ونصف دائرة على ارتفاع المكعب، فنلقي من نقطة التقاطع القطع المكافئ والدائرة خطاً إلى المحور، فذلك الخط هو الجذر المطلوب.}

\vspace{0.3cm}

\textit{لنضرب مثالاً: كعب وجذور يساويان عدداً. نجعل العدد متوازياً مستطيلات قاعدته مربعة، أضلاعها مساوية لعدد الجذور، وارتفاعه خط معلوم. ثم نأخذ قطعاً مكافئاً رأسه نقطة، ومحوره خط مستقيم، ومعلمه مساوٍ لضلع القاعدة؛ ونرسم دائرة على ارتفاع المتوازي. من نقطة التقاطع القطع المكافئ والدائرة نلقي خطاً إلى المحور: فذلك الخط هو الجذر المطلوب.}

\end{propbox}

% ============================================================================
% التركيب
% ============================================================================
\subsubsection{\ar{التركيب الهندسي}}

\ar{يتم التركيب على النحو التالي:}

\begin{enumerate}[noitemsep]
\item \ar{ليكن} $b^2 = c$ \ar{(معامل} $x$\ar{)، وليكن} $h$ \ar{بحيث} $b^2 h = d$ \ar{(الحد الثابت)}
\item \ar{نرسم \textbf{قطعاً مكافئاً} برأس} $B$ \ar{ومحور} $BZ$ \ar{ومعلم} $b$\ar{، بحيث لأي نقطة على القطع المكافئ:} $y^2 = bx$
\item \ar{ننصب عموداً} $h$ \ar{عند} $B$\ar{، وعلى} $h$ \ar{كقطر نرسم \textbf{نصف دائرة}}
\item \ar{لتقطع نصف الدائرة القطع المكافئ عند النقطة} $D$
\item \ar{من} $D$ \ar{نلقي} $DE$ \ar{عموداً على} $h$\ar{، و} $DZ$ \ar{عموداً على} $BZ$
\item \ar{فيكون} $y = BE$ \ar{هو الجذر المطلوب}
\end{enumerate}

\begin{figure}[h]
\centering
\begin{tikzpicture}[scale=1.2,>=Stealth]
    % Axes
    \draw[->] (-0.3,0) -- (5,0) node[right] {$x$};
    \draw[->] (0,-0.3) -- (0,4.5) node[above] {$y$};
    % Parabola y^2 = 2x (parameter b=2)
    \draw[thick,domain=0:4,samples=100,variable=\x,blue]
        plot ({\x},{sqrt(2*\x)}) node[right] {\small \ar{قطع مكافئ}};
    \draw[thick,domain=0:4,samples=100,variable=\x,blue]
        plot ({\x},{-sqrt(2*\x)});
    % Semicircle on height h=3 at B
    \draw[thick,red] (0,0) arc (-90:90:1.5) node[above right] {\small \ar{نصف دائرة}};
    % Intersection point D
    \coordinate (D) at (1.35, 1.65);
    \fill[black] (D) circle (1.5pt) node[above right] {$D$};
    % Drop perpendiculars
    \draw[dashed] (D) -- (1.35,0) node[below] {$Z$};
    \draw[dashed] (D) -- (0,1.65) node[left] {$E$};
    % Labels
    \node[below left] at (0,0) {$B$};
    \draw[<->] (5.3,0) -- (5.3,1.65) node[midway,right] {$y$};
    \draw[<->] (0,-0.6) -- (1.35,-0.6) node[midway,below] {$x$};
    % Parameter
    \node at (3.5,3.5) {$y^2 = bx$};
    \node at (2.5,0.8) {$x^2 = y(h-y)$};
\end{tikzpicture}
\caption{\ar{تركيب الخيام لـ} $x^3 + cx = d$ \ar{بتقاطع القطع المكافئ ونصف الدائرة}}
\label{fig:khayyam1}
\end{figure}

% ============================================================================
% البرهان
% ============================================================================
\subsubsection{\ar{البرهان}}

\ar{بخاصية القطع المكافئ:}
\[y^2 = bx \quad \Longrightarrow \quad \frac{x}{y} = \frac{y}{b}\]

\ar{بخاصية الدائرة (ذات القطر} $h$\ar{):}
\[x^2 = y(h - y)\]

\ar{من هذين:}
\[\frac{y}{b} = \frac{x}{y} = \frac{h - y}{x}\]

\ar{وبالتالي} $b : y = y : x = x : (h - y)$ \ar{في تناسب متسلسل.}

\ar{من} $b^2 : y^2 = b : (h - y)$ \ar{نحصل على:}
\[b^2(h - y) = y^3\]
\[b^2 h - b^2 y = y^3\]
\[y^3 + b^2 y = b^2 h\]

\ar{باستبدال} $c = b^2$ \ar{و} $d = b^2 h$\ar{:}
\[\boxed{y^3 + cy = d}\]

\clearpage

% ============================================================================
% الأصناف الباقية
% ============================================================================
\section{\ar{الأصناف الباقية غير القابلة للاختزال}}

% ============================================================================
% ملخص أزواج القطوع المخروطية
% ============================================================================
\subsection{\ar{ملخص أزواج القطوع المخروطية لكل صنف}}

\ar{يلخص الجدول التالي أزواج القطوع المخروطية المستخدمة لكل من الأصناف الأربعة عشر، منظمة حسب العائلات الهندسية الخمس التي وضعها الخيام:}

\begin{center}
\renewcommand{\arraystretch}{1.4}
\begin{longtable}{p{4.5cm}p{3cm}p{4.5cm}}
\toprule
\textbf{\ar{المعادلة}} & \textbf{\ar{القطوع المخروطية}} & \textbf{\ar{التسمية العربية عند الخيام}} \\
\midrule
\endfirsthead
\toprule
\textbf{\ar{المعادلة}} & \textbf{\ar{القطوع المخروطية}} & \textbf{\ar{التسمية العربية}} \\
\midrule
\endhead
\bottomrule
\endfoot
$x^3 + bx^2 = d$ & \ar{دائرة + مكافئ} & \ar{دائرة ومكافئ} \\
$x^3 + d = bx^2$ & \ar{مكافئ + زائد} & \ar{مكافئ وزائد} \\
$x^3 = bx^2 + d$ & \ar{مكافئ + زائد} & \ar{مكافئ وزائد} \\
$x^3 + cx = d$ & \ar{دائرة + مكافئ} & \ar{دائرة ومكافئ} \\
$x^3 + d = cx$ & \ar{مكافئ + زائد} & \ar{مكافئ وزائد} \\
$x^3 = cx + d$ & \ar{مكافئ + زائد} & \ar{مكافئ وزائد} \\
$x^3 + bx^2 = cx + d$ & \ar{دائرة + زائد} & \ar{دائرة وزائد} \\
$x^3 + cx = bx^2 + d$ & \ar{دائرة + زائد} & \ar{دائرة وزائد} \\
$x^3 + d = bx^2 + cx$ & \ar{دائرة + زائد} & \ar{دائرة وزائد} \\
$x^3 = bx^2 + cx + d$ & \ar{زائد + زائد} & \ar{زائدان} \\
\bottomrule
\end{longtable}
\end{center}

% ============================================================================
% في وجود الحلول وتعددها
% ============================================================================
\subsection{\ar{في وجود الحلول وتعددها}}

\ar{يناقش الخيام بعناية الشروط التي توجد عندها حلول. يلاحظ أن بعض المعادلات التكعيبية قد يكون لها:}

\begin{itemize}[noitemsep]
\item \ar{لا حل (حقيقي موجب)}
\item \ar{حل واحد موجب}
\item \ar{حلان موجبان}
\end{itemize}

\begin{propbox}[AutoFakeBold=true,AutoFakeSlant=true]
\ar{واعلم أن بعض هذه الأصناف قد يكون له حل واحد، وبعضها قد يكون له حلان، وبعضها لا يمكن أن يوجد إلا بشرط. فإذا كان القطعان لا يلتقيان فلا حل للمسألة، وإذا كانا يلتقيان في نقطة واحدة فحل واحد، وإذا كانا يلتقيان في نقطتين فحلان.}

\vspace{0.3cm}

\textit{اعلم أن بعض هذه الأصناف قد يكون له حل واحد، وبعضها قد يكون له حلان، وبعضها لا يمكن إيجاده إلا بشرط. إذا لم يتقاطع القطعان فلا حل للمسألة، وإذا تقاطعا في نقطة واحدة فحل واحد، وإذا تقاطعا في نقطتين فحلان.}
\end{propbox}

% ============================================================================
% الصنف الخامس والعشرون
% ============================================================================
\subsection{\ar{الصنف الخامس والعشرون: كعب يساوي أموالاً وجذوراً وأعداداً}}

\ar{الحالة الأخيرة ($x^3 = bx^2 + cx + d$) هي الأكثر تعقيداً:}

\begin{propbox}[AutoFakeBold=true,AutoFakeSlant=true]
\ar{وأما الصنف الخامس والعشرون: كعب يساوي أموالاً وجذوراً وأعداداً، فإنه يحتاج إلى قطعين زائدين مختلفين. فنصنع قطعاً زائداً له أحد محوريه على امتداد الخط المعلوم، ونصنع قطعاً زائداً آخر له أحد محوريه على امتداد الخط الآخر، ونلقي من نقطة التقاطعهما خطاً إلى الموضع المعلوم، فهو الجذر المطلوب.}

\vspace{0.3cm}

\textit{أما الصنف الخامس والعشرون: الكعب يساوي أموالاً وجذوراً وأعداداً، فيتطلب قطعين زائدين مختلفين. نصنع قطعاً زائداً أحد محوريه على امتداد الخط المعروف، ونصنع قطعاً زائداً آخر أحد محوريه على امتداد الخط الآخر؛ ومن نقطة تقاطعهما نلقي خطاً إلى الموضع المعروف، فهو الجذر المطلوب.}
\end{propbox}

\clearpage

% ============================================================================
% مسائل وتطبيقات
% ============================================================================
\section{\ar{مسائل وتطبيقات}}

% ============================================================================
% قسمة عشرة إلى جزأين
% ============================================================================
\subsection{\ar{قسمة عشرة إلى جزأين}}

\ar{يوضح الخيام طريقته بمسألة تؤدي إلى معادلة تكعيبية:}

\begin{propbox}[\arBold{مسألة}]
\ar{مسألة: قال قائل: قسّم عشرة إلى جزأين، بحيث يكون مجموع مربعي الجزأين مع خارج قسمة الكبير على الصغير اثنين وسبعين.}

\vspace{0.3cm}

\textit{قال شخص: اقسم عشرة إلى جزأين بحيث يكون مجموع مربعي الجزأين مع خارج قسمة الأكبر على الأصغر اثنين وسبعين.}
\end{propbox}

{\arBold{الحل.}} \ar{لنكن الجزء الأكبر} $x$\ar{. فالجزء الأصغر هو} $10 - x$\ar{. الشرط يعطي:}
\[x^2 + (10 - x)^2 + \frac{x}{10 - x} = 72\]

\ar{بتبسيط:}
\[2x^2 - 20x + 100 + \frac{x}{10 - x} = 72\]

\ar{وهذا يؤدي إلى:}
\[x^3 + 200x = 20x^2 + 2000\]

\ar{وهي من الصنف} $x^3 + cx = bx^2 + d$\ar{، وتحل بتقاطع دائرة وزائد.}

% ============================================================================
% في مضاعفة المكعب
% ============================================================================
\subsection{\ar{في مضاعفة المكعب}}

\ar{يُبين الخيام أن مسألة مضاعفة المكعب القديمة تُعادل حل المعادلة $x^3 = 2a^3$، وهي ضمن تصنيفه بصنف ``كعب يساوي عدداً'':}

\begin{propbox}[AutoFakeBold=true,AutoFakeSlant=true]
\ar{ومن ذلك مضاعفة المكعب، وهي مسألة مشهورة بين القدماء. فإنها تؤول إلى أن كعباً يساوي عدداً، والعدد هو مضاعفة المكعب المعطى.}

\vspace{0.3cm}

\textitومنها مضاعفة المكعب، وهي مسألة مشهورة بين القدماء. فإنها تختزل إلى أن الكعب يساوي عدداً، والعدد هو مضاعفة المكعب المعطى.}
\end{propbox}

\clearpage

% ============================================================================
% كلمات ختامية
% ============================================================================
\section{\ar{كلمات ختامية}}

\ar{يختتم الخيام رسالته بتأملات في طبيعة الجبر وعلاقته بالهندسة:}

\begin{propbox}[AutoFakeBold=true,AutoFakeSlant=true]
\ar{وقد أتممنا ما أردنا من جمع أصناف المسائل الجبرية، ووضع البراهين عليها. ونحن نعلم أن بعض الناس يعتقدون أن الجبر بريء من البرهان الهندسي، ولكن هذا باطل. فإن الأصناف التي ذكرناها لا يمكن أن تُبرهن إلا بالقطوع المخروطية أو بغيرها من الأدلة الهندسية. والجبر والمقابلة إنما هما وسيلة لاستخراج المجهول، والبرهان على صحة ذلك إنما هو بالهندسة.}

\vspace{0.3cm}

\textit{لقد أكملنا ما قصدناه من جمع أصناف المسائل الجبرية ووضع البراهين عليها. نعلم أن بعض الناس يعتقدون أن الجبر خالٍ من البرهان الهندسي، ولكن هذا باطل. فإن الأصناف التي ذكرناها لا يمكن إثباتها إلا بالقطوع المخروطية أو بوسائل هندسية أخرى. الجبر والمقابلة إنما هما وسيلة لاستخراج المجهول، والبرهان على صحة ذلك هو بالهندسة.}
\end{propbox}

\vspace{0.5cm}

\begin{center}
{\arLarge{والله ولي التوفيق.}}

\vspace{0.3cm}

\textit{والله مصدر كل توفيق.}
\end{center}

\vspace{1cm}

\begin{center}
{\color{rulecolor}\rule{0.6\textwidth}{0.8pt}}

\vspace{0.5cm}

{\arLarge{تمت الرسالة}}

\vspace{0.3cm}

{\small\textit{End of the Treatise}}
\end{center}

\clearpage

% ============================================================================
% ملاحظة المحرر
% ============================================================================
\section*{\ar{ملاحظة المحرر}}
\addcontentsline{toc}{section}{\ar{ملاحظة المحرر}}

\ar{هذه النسخة مبنية أساساً على مخطوطة لاهور من القرن الثالث عشر (Smith Oriental MS 45، مكتبة الكتب النادرة بجامعة كولومبيا)، مع الرجوع إلى المصادر التالية:}

\begin{itemize}[noitemsep]
\item Woepcke, F. \textit{L'Algèbre d'Omar Alkhayyâmî}. Paris, 1851.
\item Kasir, D. S. \textit{The Algebra of Omar Khayyam}. New York: Teachers College Press, 1931.
\item Rashed, R. and Vahabzadeh, B. \textit{Omar Khayyam, the Mathematician}. New York: Bibliotheca Persica Press, 2000.
\end{itemize}

\ar{تم تسوية النص العربي لتسهيل القراءة مع الحفاظ على المفردات التقنية والمحتوى الرياضي للأصل.}

\ar{تمثل التراكيب الهندسية بالرموز الحديثة مع الحفاظ على البنية المنطقية لبراهين الخيام.}

\vspace{1cm}

\begin{center}
{\color{rulecolor}\rule{0.4\textwidth}{0.4pt}}

\vspace{0.5cm}

{\footnotesize\textit{Typeset in \LaTeX\ using XeLaTeX, polyglossia, and Noto Naskh Arabic}}

{\footnotesize\textit{Mathematical typesetting via amsmath and TikZ}}
\end{center}

\end{document}
% ============================================================================
% End of Document
% ============================================================================
